\subsection{\texorpdfstring{\(R^{n}\) 的豪斯多夫商群}{R 的 n 次幂的豪斯多夫商群}}

\(\mathbf{R}^n\) 的每个豪斯多夫商群都形如 \(\mathbf{R}^n/\mathbf{H}\) ，其中 \(\mathbf{H}\) 是 \(\mathbf{R}^n\) 的闭子群（第 III 章 § 2 第 6 小节命题 18）。由第 3 小节定理 2 的推论 2，存在一个自同构 \(f\) （ \(\mathbf{R}^n\) 的），它把 \(\mathbf{H}\) 变成子群 \(\mathbf{H}'\) ，即由 \(p\) 个向量 \(\boldsymbol{e}_i\) （典范基的）所生成的向量子空间与由 \(q\) 个其余向量（共 \(n-p\) 个）所生成的离散群的直和

\[
\boldsymbol {e} _ {i}, \quad \mathrm{o} \leqslant p + q \leqslant n.
\]

过渡到商后， \(f\) 诱导出 \(\mathbf{R}^n / \mathbf{H}\) 到 \(\mathbf{R}^n / \mathbf{H}'\) 上的一个同构（第 III 章 § 2 第 8 小节注 3）；而 \(\mathbf{R}^n / \mathbf{H}'\) 同构于 \(\mathbf{R}^{n - p - q} \times \mathbf{T}^q\) （第 III 章 § 2 第 9 小节命题 26 的推论）。因此：

命题 9. \(\mathbf{R}^{n}\) 的每个豪斯多夫商群都同构于一个乘积群 \(\mathbf{R}^{h} \times \mathbf{T}^{k}\) （ \(0 \leqslant h + k \leqslant n\) ）。

乘积空间 \(\mathbf{T}^n\) （并且，沿用习惯的说法，拓扑群 \(\mathbf{T}^n\) ）称为 \(n\) 维环面；由第 V 章 § 1 第 2 小节命题 4， \(\mathbf{T}^n\) 是紧的、连通的且局部连通的。

此外，若用 C 表示 \(R^{n}\) 中边长为 1 的闭立方体，则 \(T^{n}\) 同胚于 C 关于等价关系`` \(x_{i} \equiv y_{i} \pmod{1}\) （ \(1 \leqslant i \leqslant n\) ）''（在 C 的点 \(\boldsymbol{x} = (\boldsymbol{x}_{i})\) 与 \(\boldsymbol{y} = (\boldsymbol{y}_{i})\) 之间）的商空间。于是 \(T^{n}\) 是由立方体 C 经过“把相对的面粘合起来”而得到的。

命题 10. 拓扑群 \(\mathbf{T}^n\) 局部同构于 \(\mathbf{R}^n\) 。因为 \(\mathbf{T}^n = (\mathbf{R}/\mathbf{Z})^n\) 同构于 \(\mathbf{R}^n/\mathbf{Z}^n\) （第 III 章 § 2 第 9 小节命题 26 的推论），而 \(\mathbf{Z}^n\) 是 \(\mathbf{R}^n\) 的离散子群（第 III 章 § 2 第 6 小节命题 19）。

由此得出，诸群 \(\mathbf{R}^p\times \mathbf{T}^{n - p}\) 都局部同构于 \(\mathbf{R}^n\) \((\mathrm{o}\leqslant p\leqslant n)\) ；在 § 2 第 2 小节中我们将看到，它们是具有此性质的仅有的连通群。

